% Part: first-order-logic
% Chapter: model-theory
% Section: theory-of-m

\documentclass[../../../include/open-logic-section]{subfiles}

\begin{document}

\olfileid[pt]{mod}{bas}{thm}

\section{A Teoria de uma \printtoken{S}{structure}}

Cada !!{structure}~$\Struct{M}$ torna algumas !!{sentence}s verdadeiras e
outras falsas. O conjunto de todas as !!{sentence}s que ela torna verdadeiras é chamado
de sua \emph{teoria}. Esse conjunto é de fato uma teoria, pois tudo o que ele
implica deve ser verdadeiro em todos os seus modelos, inclusive~$\Struct{M}$.

\begin{defn}
  Dada !!a{structure}~$\Struct M$, a \emph{teoria} de
  $\Struct{M}$ é o conjunto $\Theory{M}$ das !!{sentence}s
  verdadeiras em $\Struct{M}$, O.s.\@, $\Theory{M} =
  \Setabs{!A}{\Sat{M}{!A}}$.
\end{defn}

Também usamos informalmente o termo ``teoria'' para nos referir a conjuntos
de !!{sentence}s que têm uma interpretação pretendida, sejam eles dedutivamente
fechados ou não.

\begin{prop}
Para toda $\Struct{M}$, $\Theory{M}$ é completa.
\end{prop}

\begin{proof}
Para toda !!{sentence}~$!A$, ou $\Sat{M}{!A}$ ou $\Sat{M}{\lnot !A}$;
portanto, ou $!A \in \Theory{M}$ ou $\lnot !A \in \Theory{M}$.
\end{proof}

\begin{prop}\ollabel{prop:equiv}
  Se $\Struct{N} \models !A$ para todo $!A \in \Theory{M}$, então
  $\Struct{M} \elemequiv \Struct{N}$.
\end{prop}

\begin{proof}
Como $\Sat{N}{!A}$ para todo $!A \in \Theory{M}$, $\Theory{M} \subseteq
\Theory{N}$. Se $\Sat{N}{!A}$, então $\Sat/{N}{\lnot !A}$, logo $\lnot !A
\notin \Theory{M}$. Como $\Theory{M}$ é completa, $!A \in
\Theory{M}$. Portanto, $\Theory{N} \subseteq \Theory{M}$, e temos
$\Struct{M} \elemequiv \Struct{N}$.
\end{proof}

\begin{rem}\ollabel{remark:R}
  Considere $\Struct{R} = \langle\Real, <\rangle$, a !!{structure}
  cujo domínio é o conjunto $\Real$ dos números reais, na !!{language}
  que contém apenas um !!{predicate} binário interpretado como a relação $<$
  sobre os reais. Claramente, $\Struct{R}$ é !!{nonenumerable};
  entretanto, como $\Theory{R}$ é obviamente consistente, pelo teorema de
  L\"owenheim--Skolem ela tem !!a{enumerable} modelo, digamos
  $\Struct{S}$, e, por \olref{prop:equiv}, $\Struct{R}
  \equiv \Struct{S}$. Além disso, como $\Struct{R}$ e $\Struct{S}$
  não são isomorfas, isso mostra que a recíproca de
  \olref[iso]{thm:isom} falha em geral.
\end{rem}

\end{document}
